\documentclass[10pt]{article}
\usepackage[a4paper,margin=0.55in]{geometry}
\usepackage{amsmath,amssymb,mathtools}
\usepackage{xcolor}
\usepackage{hyperref}
\usepackage{url}
\setlength{\jot}{5pt}
\numberwithin{equation}{section}
\allowdisplaybreaks

\newcommand{\Cz}{\mathcal{C}^{(b)}_{z}}
\newcommand{\UVb}{\left[\frac{2}{\epsilon}+2\gamma_{E}+\log\frac{(y-z)^{2}\mu^{2}_{\overline{MS}}}{4}\right]}
\newcommand{\Wl}{\Big[U^{ab'}(x')-U^{ab'}(w')\Big]}

\title{$N_{2}$ of Group I, Row III: UV logarithms, finite terms and the longitudinal logarithms, step by step}
\date{}

\begin{document}
\maketitle

\section{Starting expression}
\begin{align}
    \frac{d^{3}N_{2\,\rm NLO}^{\rm}}{d^{3}k}  & =  \frac{1}{(2\pi)^{3}}\, \frac{g^{4}}{8\pi^{4}}\,\frac{1}{k^{+}}\,\int_{\Lambda}^{k^{+}-\Lambda}\, \frac{dp^{+}}{2\pi} \, \int_{x,y,z,x',w'}\,e^{-ik\cdot(w'-z)}\,e^{-ik\cdot (z-x)\,\frac{p^{+}}{k^{+}}}\,\frac{1}{\left[p^{+}\,(y-x)^{2}+(k^{+}-p^{+})\,(y-z)^{2}\right]} \notag \\ & \times\frac{(x'-w')^{i}\,(z-x)^{m}}{(x'-w')^{2}\,(z-x)^{2}}\Bigg[\frac{\delta^{im}\,d_{\perp}}{2\,(x-z)^{2}}\,\left[(y-x)^{2}-(y-z)^{2}\right]\,\frac{p^{+}\,(k^{+}-p^{+})}{k^{+\,2}}\notag\\ &-\frac{p^{+}}{(k^{+}-p^{+})}\,\left[\frac{(y-x)^{i}\,(x-z)^{m}}{(x-z)^{2}}-\frac{(y-x)^{i}\,(y-z)^{m}}{2\,(y-z)^{2}}\right]\notag \\ & +\frac{(k^{+}-p^{+})\,\delta_{im}}{k^{+}}\,\left[\frac{(y-z)\cdot (x-z)}{(x-z)^{2}}+\frac{(y-x)\cdot (y-z)}{2\,(y-x)^{2}}-\frac{\left[(y-x)^{2}-(y-z)^{2}\right]}{2\,(x-z)^{2}}\right]\notag \\ & -\frac{(k^{+}-p^{+})}{p^{+}}\,\left[\frac{(y-z)^{i}\,(x-z)^{m}}{(x-z)^{2}}+\frac{(y-x)^{m}\,(y-z)^{i}}{2\,(y-x)^{2}}\right]\notag \\ & +\frac{p^{+}\,\delta_{im}}{k^{+}}\,\left[\frac{(y-x)\cdot (x-z)}{(x-z)^{2}}-\frac{(y-x)\cdot (y-z)}{2\,(y-z)^{2}}-\frac{\left[(y-x)^{2}-(y-z)^{2}\right]}{2\,(x-z)^{2}}\right]\notag \\ &-\left[\frac{(y-x)^{m}\,(x-z)^{i}}{(x-z)^{2}}-\frac{(y-x)^{m}\,(y-z)^{i}}{2\,(y-z)^{2}}+\frac{(y-z)^{m}\,(x-z)^{i}}{(x-z)^{2}}+\frac{(y-x)^{i}\,(y-z)^{m}}{2\,(y-x)^{2}}\right]\Bigg]\notag \\ & \times\Bigg[U^{ab'}(x')-U^{ab'}(w')\Bigg]\notag\\ &\times\Bigg[-f^{adc}\,f^{fbe}\,U^{db}(x)\,U^{ce}(z)+N_{c}\,U^{af}(y)+N_{c}\,U^{af}(x) -N_{c}\,U^{af}(x)\Bigg]\,\rho^{b'}(x')\,\rho^{f}(y)
\label{eq:N2}
\end{align}

\section{Step 0: shorthand}
\begin{align}
&\mathcal{N}\equiv\frac{1}{(2\pi)^{3}}\frac{g^{4}}{8\pi^{4}}\frac{1}{k^{+}},\qquad \frac{\mathcal{N}}{2\pi}=\frac{1}{(2\pi)^{3}}\frac{g^{4}}{16\pi^{5}}\frac{1}{k^{+}},\\ &\xi\equiv\frac{p^{+}}{k^{+}},\qquad \lambda\equiv\frac{\Lambda}{k^{+}},\\
&\Phi(p^{+})=e^{-ik\cdot(z-x)\frac{p^{+}}{k^{+}}},\qquad D(p^{+})=p^{+}(y-x)^{2}+(k^{+}-p^{+})(y-z)^{2},\\ &T^{im}=\frac{(x'-w')^{i}(z-x)^{m}}{(x'-w')^{2}(z-x)^{2}},\\
&B^{im}=\frac{(y-z)^{i}(x-z)^{m}}{(x-z)^{2}}+\frac{(y-x)^{m}(y-z)^{i}}{2(y-x)^{2}},\\
&A^{im}=\frac{(y-x)^{i}(x-z)^{m}}{(x-z)^{2}}-\frac{(y-x)^{i}(y-z)^{m}}{2(y-z)^{2}},
\end{align}
and $R^{im}(\xi)$ denotes the four non-singular brackets of \eqref{eq:N2} (the $d_{\perp}$, $(k^{+}-p^{+})\delta$, $p^{+}\delta$ and last brackets). In this notation the bracket of \eqref{eq:N2} is
\begin{equation}
R^{im}(\xi)-\frac{\xi}{1-\xi}A^{im}-\frac{1-\xi}{\xi}B^{im}.
\end{equation}
Split the colour factor into the $N_{2a}$ and $N_{2b}$ parts:
\begin{align}
\mathcal{C}^{(a)}&=\Wl\Big[-f^{adc}f^{fbe}U^{db}(x)U^{ce}(z)+N_{c}U^{af}(x)\Big]\rho^{b'}(x')\rho^{f}(y),\\
\mathcal{C}^{(b)}&=\Wl\Big[N_{c}U^{af}(y)-N_{c}U^{af}(x)\Big]\rho^{b'}(x')\rho^{f}(y),\\
\Cz&=\Wl\Big[N_{c}U^{af}(y)-N_{c}U^{af}(z)\Big]\rho^{b'}(x')\rho^{f}(y),\\
\mathcal{C}_{2}&=\mathcal{C}^{(a)}+\mathcal{C}^{(b)}=\Wl\Big[-f^{adc}f^{fbe}U^{db}(x)U^{ce}(z)+N_{c}U^{af}(y)\Big]\rho^{b'}(x')\rho^{f}(y),
\end{align}
and define
\begin{equation}
K_{S}=\frac{(x'-w')\cdot(y-z)}{(x'-w')^{2}(y-z)^{2}}\,\frac{(y-x)\cdot(y-z)}{(y-x)^{2}(x-z)^{2}}.
\end{equation}

\section{Step 1: where the $x$ integral can diverge}
Fix $0<p^{+}<k^{+}$; the only variable that is integrated over and appears in the kernel singularly is $x$. Check the three regions.
\begin{itemize}
\item \textbf{Large $x$.} $T\sim1/|x|$, $D\sim p^{+}x^{2}$, and all brackets are at most $O(1)$, so the integrand falls at least like $1/|x|^{3}$. This is finite.
\item \textbf{$x\to y$.} The brackets have at most $1/|y-x|$, which is integrable in two dimensions. At $p^{+}\to k^{+}$ one has $D\to k^{+}(y-x)^{2}$, but $\int_{0}^{k^{+}}dp^{+}/D=\frac{\log[(y-z)^{2}/(y-x)^{2}]}{(y-z)^{2}-(y-x)^{2}}$, which is only a logarithm and is integrable. This is finite.
\item \textbf{$x\to z$.} Several terms go as $1/(x-z)^{2}$, which is logarithmically divergent in two dimensions. This is the UV region; Step 2 works it out.
\end{itemize}

\section{Step 2: the integrand as $x\to z$}
Put $r=x-z\to0$, $n=y-z$, $a=x'-w'$, so $y-x=n-r$ and $(y-x)^{2}-(y-z)^{2}=-2\,n\cdot r+r^{2}$. At $r=0$: $\Phi=1$ and $D=k^{+}n^{2}$. Keep only the $1/r^{2}$ terms and average over the direction of $r$ with $\langle r^{i}r^{j}\rangle=\frac{r^{2}}{2}\delta^{ij}$ (in $d$ dimensions $\frac{r^{2}}{d}\delta^{ij}$ together with $d_{\perp}=d$, which gives the same results).
\begin{align}
\text{$d_{\perp}$ term:}\quad& \frac{d_{\perp}}{2}\,\xi(1-\xi)\,\frac{(-a\cdot r)(-2\,n\cdot r)}{a^{2}r^{4}}\ \to\ \xi(1-\xi)\,\frac{a\cdot n}{a^{2}r^{2}},\\
\text{$(1-\xi)\delta$ term:}\quad& \text{bracket}=\frac{2\,n\cdot r}{r^{2}}+O(1)\ \Rightarrow\ (1-\xi)\frac{(-a\cdot r)(2\,n\cdot r)}{a^{2}r^{4}}\ \to\ -(1-\xi)\frac{a\cdot n}{a^{2}r^{2}},\\
\text{$-\tfrac{1-\xi}{\xi}B$ term:}\quad& T^{im}B^{im}=-\frac{a\cdot n}{a^{2}r^{2}}+O(1/r)\ \Rightarrow\ +\frac{1-\xi}{\xi}\frac{a\cdot n}{a^{2}r^{2}},\\
\text{$\xi\delta$ term:}\quad& \text{bracket}=\frac{2\,n\cdot r}{r^{2}}+O(1)\ \Rightarrow\ -\xi\,\frac{a\cdot n}{a^{2}r^{2}},\\
\text{last bracket:}\quad& -\frac{2\,n^{m}r^{i}}{r^{2}}+O(1)\ \Rightarrow\ \frac{a^{i}(-r^{m})}{a^{2}r^{2}}\Big(-\frac{2n^{m}r^{i}}{r^{2}}\Big)\to+\frac{a\cdot n}{a^{2}r^{2}},\\
\text{$-\tfrac{\xi}{1-\xi}A$ term:}\quad& T^{im}A^{im}=-\frac{a\cdot n}{a^{2}r^{2}}+O(1/r)\ \Rightarrow\ +\frac{\xi}{1-\xi}\frac{a\cdot n}{a^{2}r^{2}}.
\end{align}
Adding these up, the constant terms cancel ($-(1-\xi)-\xi+1=0$), and with $1/D\to1/(k^{+}n^{2})$:
\begin{equation}
\text{integrand of \eqref{eq:N2}}\ \xrightarrow{x\to z}\ \frac{1}{k^{+}}\,P(\xi)\,\frac{(x'-w')\cdot(y-z)}{(x'-w')^{2}(y-z)^{2}\,(x-z)^{2}}\times(\text{colour at }x=z),\qquad P(\xi)=\xi(1-\xi)+\frac{\xi}{1-\xi}+\frac{1-\xi}{\xi}.
\label{eq:limit}
\end{equation}
This was also checked numerically: the angle-averaged integrand divided by $K_{S}$ tends to $P(\xi)$ (for example $4.41$, $2.25$ and $9.2011$ at $\xi=0.2,\,0.5,\,0.9$), and the difference falls like $r^{2}$.

\section{Step 3: colour at $x\to z$, and why only $N_{2b}$ has a UV divergence}
Adjoint Wilson lines leave $f$ invariant, so $U^{db}(z)U^{ce}(z)f^{fbe}=f^{dcg}U^{gf}(z)$. Then
\begin{equation}
-f^{adc}f^{fbe}U^{db}(z)U^{ce}(z)+N_{c}U^{af}(z)=-f^{adc}f^{dcg}U^{gf}(z)+N_{c}U^{af}(z)=-N_{c}U^{af}(z)+N_{c}U^{af}(z)=0 .
\end{equation}
This was checked numerically for a random SU(3) matrix: the residual is $2\times10^{-15}$.
\begin{itemize}
\item \textbf{$N_{2a}$} (colour $\mathcal{C}^{(a)}$): the colour vanishes linearly at $x\to z$, so the integrand is at most $O(1/r)$. $N_{2a}$ is \textbf{UV finite}.
\item \textbf{$N_{2b}$} (colour $\mathcal{C}^{(b)}$): $\mathcal{C}^{(b)}\to\Cz\neq0$, so $N_{2b}$ is \textbf{UV divergent}, with the coefficient given by \eqref{eq:limit}.
\end{itemize}

\section{Step 4: subtraction}
Define the subtraction term, which has exactly the $x\to z$ behaviour \eqref{eq:limit}:
\begin{equation}
S(p^{+})=\frac{1}{k^{+}}P(\xi)\,K_{S}\,\Cz .
\end{equation}
At $x\to z$ we have $\frac{(y-x)\cdot(y-z)}{(y-x)^{2}}\to1$, so $K_{S}\to\frac{(x'-w')\cdot(y-z)}{(x'-w')^{2}(y-z)^{2}(x-z)^{2}}$. In addition $K_{S}$ is integrable at $x\to y$ and at large $x$ (where it falls like $|x|^{-3}$), so $S$ adds no new divergence. Write, identically,
\begin{equation}
N_{2}=\underbrace{\mathcal{N}\!\int\!\frac{dp^{+}}{2\pi}\!\int e^{-ik\cdot(w'-z)}\,S}_{N_{2b\,\rm UV}}\;+\;\underbrace{\mathcal{N}\!\int\!\frac{dp^{+}}{2\pi}\!\int e^{-ik\cdot(w'-z)}\Big[\text{integrand of }\eqref{eq:N2}-S\Big]}_{N_{2\,\rm finite}}.
\end{equation}
The second term is UV finite, so we set $d=2$ ($d_{\perp}=2$) in it. All the UV divergence is in the first term.

\section{Step 5: UV part}
\subsection*{5a. The $x$ integral in two dimensions (to see the divergence)}
Put $x=z+t$, $s=|t|$, $n=y-z$. Writing $n^{2}-ns\cos\varphi=\tfrac12\big[(n^{2}+s^{2}-2ns\cos\varphi)+(n^{2}-s^{2})\big]$ and using $\int_{0}^{2\pi}\frac{d\varphi}{A-B\cos\varphi}=\frac{2\pi}{\sqrt{A^{2}-B^{2}}}$:
\begin{equation}
\int_{0}^{2\pi}d\varphi\,\frac{n^{2}-ns\cos\varphi}{n^{2}+s^{2}-2ns\cos\varphi}=\pi+\pi\,\frac{n^{2}-s^{2}}{|n^{2}-s^{2}|}=\begin{cases}2\pi,&s<|n|\\0,&s>|n|\end{cases}
\end{equation}
(checked numerically). Hence
\begin{equation}
\int_{x}\frac{(y-x)\cdot(y-z)}{(y-x)^{2}(x-z)^{2}}=2\pi\int_{0}^{|y-z|}\frac{ds}{s},
\end{equation}
which is logarithmically divergent at $s\to0$, that is at $x\to z$ (UV). It is finite at large $x$ and at $x\to y$.

\subsection*{5b. The same integral in $d=2-\epsilon$}
Using a Feynman parameter, $\frac{1}{(y-x)^{2}(x-z)^{2}}=\int_{0}^{1}\frac{d\alpha}{[(t-\alpha n)^{2}+\alpha(1-\alpha)n^{2}]^{2}}$, and shifting $t\to t+\alpha n$:
\begin{align}
\mu^{-\epsilon}\!\int d^{d}t\,\frac{(n-t)\cdot n}{(n-t)^{2}t^{2}}&=\mu^{-\epsilon}\,n^{2}\,\pi^{d/2}\Gamma\big(2-\tfrac d2\big)\int_{0}^{1}d\alpha\,(1-\alpha)\big[\alpha(1-\alpha)n^{2}\big]^{\frac d2-2}\notag\\
&=\pi^{d/2}\,\Gamma\big(2-\tfrac d2\big)\frac{\Gamma(\frac d2-1)\Gamma(\frac d2)}{\Gamma(d-1)}\,(\mu^{2}n^{2})^{-\epsilon/2}\notag\\
&=\pi\left[\frac{2}{\epsilon_{\rm UV}}+2\gamma_{E}+\log\frac{(y-z)^{2}\mu^{2}_{\overline{MS}}}{4}\right]+O(\epsilon),
\end{align}
with $\mu^{2}_{\overline{MS}}=4\pi e^{-\gamma_{E}}\mu^{2}$ and $\epsilon_{\rm UV}=-\epsilon$; this was checked with sympy. This is the same convention as the Divergent file.

\subsection*{5c. The $p^{+}$ integral, done exactly}
$K_{S}\Cz$ does not depend on $p^{+}$, so
\begin{align}
\frac{1}{k^{+}}\int_{\Lambda}^{k^{+}-\Lambda}dp^{+}P(\xi)&=\int_{\lambda}^{1-\lambda}d\xi\Big[\xi(1-\xi)+\Big(\frac{1}{1-\xi}-1\Big)+\Big(\frac{1}{\xi}-1\Big)\Big]\notag\\
&=\Big(\frac16-\lambda^{2}+\frac{2\lambda^{3}}{3}\Big)+2\log\frac{1-\lambda}{\lambda}-2(1-2\lambda)
=2\log\frac{k^{+}-\Lambda}{\Lambda}-\frac{11}{6}+\frac{4\Lambda}{k^{+}}-\frac{\Lambda^{2}}{k^{+2}}+\frac{2\Lambda^{3}}{3k^{+3}} .
\end{align}

\subsection*{5d. Result for the UV part}
The prefactor is $\frac{\mathcal{N}}{2\pi}\cdot\pi\cdot N_{c}=\frac{1}{(2\pi)^{3}}\frac{g^{4}N_{c}}{16\pi^{4}}\frac{1}{k^{+}}$, so
\begin{align}
\frac{d^{3}N_{2b\,\rm UV\, NLO}}{d^{3}k} & = \frac{1}{(2\pi)^{3}}\,\frac{g^{4}N_{c}}{16\pi^{4}}\,\frac{1}{k^{+}}
\left[2\log\Big(\frac{k^{+}-\Lambda}{\Lambda}\Big)-\frac{11}{6}+\frac{4\Lambda}{k^{+}}-\frac{\Lambda^{2}}{k^{+2}}+\frac{2\Lambda^{3}}{3k^{+3}}\right]\notag\\
&\times\int_{y,z,x',w'}e^{-ik\cdot(w'-z)}\UVb\notag\\
&\times\frac{(x'-w')\cdot(y-z)}{(x'-w')^{2}\,(y-z)^{2}}\Wl\Big[U^{af}(y)-U^{af}(z)\Big]\rho^{b'}(x')\,\rho^{f}(y)\\
&\xrightarrow{\Lambda\to0}\ \frac{1}{(2\pi)^{3}}\,\frac{g^{4}N_{c}}{16\pi^{4}}\,\frac{1}{k^{+}}\left[{\color{blue}2\log\frac{k^{+}}{\Lambda}}-\frac{11}{6}\right]\int_{y,z,x',w'}(\cdots).
\end{align}
Here $2\log\frac{k^{+}}{\Lambda}=2\log\frac{\vee}{\Lambda}-2\log\frac{\vee}{k^{+}}$. The $-\frac{11}{6}\times$UV piece is the one that combines with Row II ($\frac{11}{3}$) in the Divergent file. The blue piece is a longitudinal log multiplying the UV bracket; Step 7 shows that it cancels against the $-2\log$ term of the finite part when the $x$ integral is undone.

\section{Step 6: finite part (before the plus prescription)}
\begin{align}
\frac{d^{3}N_{2\,\rm finite}}{d^{3}k}&=\mathcal{N}\int_{\Lambda}^{k^{+}-\Lambda}\frac{dp^{+}}{2\pi}\int_{x,y,z,x',w'}e^{-ik\cdot(w'-z)}\Bigg\{\frac{\Phi}{D}\,T^{im}\Big[R^{im}(\xi)-\frac{\xi}{1-\xi}A^{im}-\frac{1-\xi}{\xi}B^{im}\Big]\mathcal{C}_{2}-\frac{1}{k^{+}}P(\xi)\,K_{S}\,\Cz\Bigg\},
\label{eq:fin0}
\end{align}
with $d_{\perp}=2$. This is UV finite at every $p^{+}$ by Steps 2 to 4, but it still has the $\Lambda$ cut-offs. The singular coefficients are $\frac{1-\xi}{\xi}$ (at $p^{+}\to0$) and $\frac{\xi}{1-\xi}$ (at $p^{+}\to k^{+}$), together with the same two inside $P(\xi)$.

\section{Step 7: plus prescription and the longitudinal logarithms}
\subsection*{7a. Definitions}
\[
\Big[\frac{k^{+}}{p^{+}}\Big]_{+}f(p^{+})\equiv\frac{k^{+}}{p^{+}}\big[f(p^{+})-f(0)\big],\qquad
\Big[\frac{k^{+}}{q^{+}}\Big]_{+}h(q^{+})\equiv\frac{k^{+}}{q^{+}}\big[h(q^{+})-h(0)\big],\]\[
\Big[\frac{k^{+}}{k^{+}-p^{+}}\Big]_{+}f(p^{+})\equiv\frac{k^{+}}{k^{+}-p^{+}}\big[f(p^{+})-f(k^{+})\big].
\]
Then, exactly,
\begin{equation}
\int_{\Lambda}^{k^{+}-\Lambda}dp^{+}\frac{k^{+}}{p^{+}}f=\int_{\Lambda}^{k^{+}-\Lambda}dp^{+}\Big[\frac{k^{+}}{p^{+}}\Big]_{+}f+k^{+}f(0)\log\frac{k^{+}-\Lambda}{\Lambda}
\ \xrightarrow{\Lambda\to0}\ \int_{0}^{k^{+}}dp^{+}\Big[\frac{k^{+}}{p^{+}}\Big]_{+}f+k^{+}f(0)\log\frac{k^{+}-\Lambda}{\Lambda}.
\end{equation}

\subsection*{7b. The $p^{+}\to0$ end of the main term}
Write $-\frac{1-\xi}{\xi}B=-\frac{k^{+}}{p^{+}}B+B$ and set $f(p^{+})=\Phi\,T^{im}B^{im}/D$. At $p^{+}=0$: $\Phi=1$ and $D=k^{+}(y-z)^{2}$, so
\begin{align}
T^{im}B^{im}&=\frac{(x'-w')\cdot(y-z)\,(x-z)\cdot(z-x)}{(x'-w')^{2}(z-x)^{2}(x-z)^{2}}+\frac{(x'-w')\cdot(y-z)\,(y-x)\cdot(z-x)}{2(x'-w')^{2}(z-x)^{2}(y-x)^{2}},\\
-k^{+}f(0)&=+\frac{(x'-w')\cdot(y-z)}{(x'-w')^{2}(y-z)^{2}}\left[\frac{1}{(x-z)^{2}}+\frac{(x-z)\cdot(y-x)}{2(x-z)^{2}(y-x)^{2}}\right].
\end{align}

\subsection*{7c. The $p^{+}\to k^{+}$ end of the main term (in $q^{+}=k^{+}-p^{+}$)}
With $p^{+}=k^{+}-q^{+}$ and $\int_{\Lambda}^{k^{+}-\Lambda}dp^{+}=\int_{\Lambda}^{k^{+}-\Lambda}dq^{+}=-\int_{k^{+}-\Lambda}^{\Lambda}dq^{+}$, the term $-\frac{\xi}{1-\xi}A=-\big(\frac{k^{+}}{q^{+}}-1\big)A$ becomes
\begin{equation}
\int_{k^{+}-\Lambda}^{\Lambda}dq^{+}\Big(\frac{k^{+}}{q^{+}}-1\Big)h(q^{+})=\int_{k^{+}}^{0}dq^{+}\Big\{\Big[\frac{k^{+}}{q^{+}}\Big]_{+}-1\Big\}h(q^{+})-k^{+}h(0)\log\frac{k^{+}-\Lambda}{\Lambda},
\end{equation}
where $h(q^{+})=e^{-ik\cdot(z-x)\frac{k^{+}-q^{+}}{k^{+}}}\,T^{im}A^{im}\big/\big[(k^{+}-q^{+})(y-x)^{2}+q^{+}(y-z)^{2}\big]$. At $q^{+}=0$: $\widetilde D=k^{+}(y-x)^{2}$ and $e^{-ik\cdot(w'-z)}e^{-ik\cdot(z-x)}=e^{-ik\cdot(w'-x)}$, so
\begin{align}
T^{im}A^{im}&=\frac{(x'-w')\cdot(y-x)\,(z-x)\cdot(x-z)}{(x'-w')^{2}(z-x)^{2}(x-z)^{2}}-\frac{(x'-w')\cdot(y-x)\,(z-x)\cdot(y-z)}{2(x'-w')^{2}(z-x)^{2}(y-z)^{2}},\\
-k^{+}h(0)&=+e^{-ik\cdot(z-x)}\frac{(x'-w')\cdot(y-x)}{(x'-w')^{2}(y-x)^{2}}\left[\frac{1}{(x-z)^{2}}-\frac{(x-z)\cdot(y-z)}{2(x-z)^{2}(y-z)^{2}}\right].
\end{align}
The endpoint coefficients of 7b and 7c were checked numerically: for random points and $\Lambda/k^{+}=10^{-7}$, the original $\int_{\Lambda}^{k^{+}-\Lambda}$ agrees with the plus form plus the log terms to $10^{-8}$, i.e.\ to $O(\Lambda/k^{+})$.

\subsection*{7d. The subtraction term}
Write $P(\xi)=\xi(1-\xi)+\frac{k^{+}}{k^{+}-p^{+}}-1+\frac{k^{+}}{p^{+}}-1$. Since $K_{S}\Cz$ is $p^{+}$ independent, each endpoint gives $K_{S}\Cz\log\frac{k^{+}-\Lambda}{\Lambda}$:
\begin{equation}
-\frac{1}{k^{+}}\int_{\Lambda}^{k^{+}-\Lambda}dp^{+}P(\xi)\,K_{S}\Cz=-\frac{1}{k^{+}}\int_{0}^{k^{+}}dp^{+}\Big[\xi(1-\xi)+\Big[\frac{k^{+}}{k^{+}-p^{+}}\Big]_{+}+\Big[\frac{k^{+}}{p^{+}}\Big]_{+}-2\Big]K_{S}\Cz-2\log\frac{k^{+}-\Lambda}{\Lambda}\,K_{S}\Cz .
\end{equation}
The first term is left unevaluated, as decided. Evaluating it would give $+\frac{11}{6}K_{S}\Cz$, because $\int[\cdot]_{+}\times\text{const}=0$.

\section{Step 8: result}
\begin{equation}
\frac{d^{3}N_{2\,\rm NLO}}{d^{3}k}=\frac{d^{3}N_{2b\,\rm UV}}{d^{3}k}\ (\text{Step 5d})+\frac{d^{3}N_{2\,\rm finite}}{d^{3}k}\ (\text{8a})+\frac{d^{3}N_{2\,\log}}{d^{3}k}\ (\text{8b})+O(\Lambda/k^{+}).
\end{equation}

\subsection*{8a. Finite part ($\Lambda$ independent; $\mathcal{C}_{2}=\mathcal{C}^{(a)}+\mathcal{C}^{(b)}$)}
\begin{align}
    \frac{d^{3}N_{2\,\rm finite\, NLO}^{\rm}}{d^{3}k} & = \frac{1}{(2\pi)^{3}}\,\frac{g^{4}}{8\pi^{4}}\,\frac{1}{k^{+}}\int_{0}^{k^{+}}\frac{dp^{+}}{2\pi}\int_{x,y,z,x',w'}e^{-ik\cdot(w'-z)}\,\Bigg\{e^{-ik\cdot(z-x)\frac{p^{+}}{k^{+}}}\,\frac{1}{\left[p^{+}(y-x)^{2}+(k^{+}-p^{+})(y-z)^{2}\right]}\notag \\ & \times\frac{(x'-w')^{i}\,(z-x)^{m}}{(x'-w')^{2}\,(z-x)^{2}}\Bigg[\frac{\delta^{im}\,d_{\perp}}{2\,(x-z)^{2}}\left[(y-x)^{2}-(y-z)^{2}\right]\frac{p^{+}(k^{+}-p^{+})}{k^{+2}}\notag\\
&+\frac{(k^{+}-p^{+})\,\delta_{im}}{k^{+}}\left[\frac{(y-z)\cdot(x-z)}{(x-z)^{2}}+\frac{(y-x)\cdot(y-z)}{2(y-x)^{2}}-\frac{(y-x)^{2}-(y-z)^{2}}{2(x-z)^{2}}\right]\notag\\
&-\left[\frac{k^{+}}{p^{+}}\right]_{+}\,\left[\frac{(y-z)^{i}(x-z)^{m}}{(x-z)^{2}}+\frac{(y-x)^{m}(y-z)^{i}}{2(y-x)^{2}}\right]+\left[\frac{(y-z)^{i}(x-z)^{m}}{(x-z)^{2}}+\frac{(y-x)^{m}(y-z)^{i}}{2(y-x)^{2}}\right]\notag\\
&+\frac{p^{+}\,\delta_{im}}{k^{+}}\left[\frac{(y-x)\cdot(x-z)}{(x-z)^{2}}-\frac{(y-x)\cdot(y-z)}{2(y-z)^{2}}-\frac{(y-x)^{2}-(y-z)^{2}}{2(x-z)^{2}}\right]\notag\\
&-\left[\frac{(y-x)^{m}(x-z)^{i}}{(x-z)^{2}}-\frac{(y-x)^{m}(y-z)^{i}}{2(y-z)^{2}}+\frac{(y-z)^{m}(x-z)^{i}}{(x-z)^{2}}+\frac{(y-x)^{i}(y-z)^{m}}{2(y-x)^{2}}\right]\Bigg]\notag\\
&\times\Wl\Big[-f^{adc}f^{fbe}U^{db}(x)U^{ce}(z)+N_{c}\,U^{af}(y)\Big]\rho^{b'}(x')\,\rho^{f}(y)\notag\\
&-\frac{1}{k^{+}}\left[\frac{p^{+}(k^{+}-p^{+})}{k^{+2}}+\left[\frac{k^{+}}{k^{+}-p^{+}}\right]_{+}+\left[\frac{k^{+}}{p^{+}}\right]_{+}-2\right]
\frac{(x'-w')\cdot(y-z)}{(x'-w')^{2}\,(y-z)^{2}}\,\frac{(y-x)\cdot(y-z)}{(y-x)^{2}\,(x-z)^{2}}\notag\\
&\times\Wl\Big[N_{c}\,U^{af}(y)-N_{c}\,U^{af}(z)\Big]\rho^{b'}(x')\,\rho^{f}(y)\Bigg\}\notag \\ & + \frac{1}{(2\pi)^{3}}\,\frac{g^{4}}{8\pi^{4}}\,\frac{1}{k^{+}}\int_{k^{+}}^{0}\frac{dq^{+}}{2\pi}\int_{x,y,z,x',w'}e^{-ik\cdot(w'-z)}\ e^{-ik\cdot(z-x)\frac{(k^{+}-q^{+})}{k^{+}}}\,\frac{1}{\left[(k^{+}-q^{+})(y-x)^{2}+q^{+}\,(y-z)^{2}\right]}\notag \\ & \times\frac{(x'-w')^{i}\,(z-x)^{m}}{(x'-w')^{2}\,(z-x)^{2}}\,\left[\frac{(y-x)^{i}(x-z)^{m}}{(x-z)^{2}}-\frac{(y-x)^{i}(y-z)^{m}}{2(y-z)^{2}}\right]\,\left\{\left[\frac{k^{+}}{q^{+}}\right]_{+}-1\right\}\notag \\ & \times\Wl\Big[-f^{adc}f^{fbe}U^{db}(x)U^{ce}(z)+N_{c}\,U^{af}(y)\Big]\rho^{b'}(x')\,\rho^{f}(y),
\label{eq:finite}
\end{align}
with $d_{\perp}=2$. To get the $N_{2a}$ and $N_{2b}$ pieces separately, replace the colour factor $\big[-ffUU+N_{c}U(y)\big]$ by $\big[-ffUU+N_{c}U(x)\big]$ for $N_{2a}$, and by $\big[N_{c}U(y)-N_{c}U(x)\big]$ for $N_{2b}$. The subtraction line belongs to $N_{2b}$ only.

\subsection*{8b. Longitudinal logarithms (these go to the JIMWLK file)}
\begin{align}
\frac{d^{3}N_{2\,\log}}{d^{3}k}&=\frac{1}{(2\pi)^{3}}\,\frac{g^{4}}{16\pi^{5}}\,\frac{1}{k^{+}}\,{\color{blue}\log\Big[\frac{k^{+}-\Lambda}{\Lambda}\Big]}\int_{x,y,z,x',w'}\Bigg\{
e^{-ik\cdot(w'-z)}\frac{(x'-w')\cdot(y-z)}{(x'-w')^{2}(y-z)^{2}}\left[\frac{1}{(x-z)^{2}}+\frac{(x-z)\cdot(y-x)}{2(x-z)^{2}(y-x)^{2}}\right]\mathcal{C}_{2}\notag\\
&\qquad+e^{-ik\cdot(w'-x)}\frac{(x'-w')\cdot(y-x)}{(x'-w')^{2}(y-x)^{2}}\left[\frac{1}{(x-z)^{2}}-\frac{(x-z)\cdot(y-z)}{2(x-z)^{2}(y-z)^{2}}\right]\mathcal{C}_{2}
-2\,e^{-ik\cdot(w'-z)}\,K_{S}\,\Cz\Bigg\},
\label{eq:logs}
\end{align}
where $\log\frac{k^{+}-\Lambda}{\Lambda}\to\log\frac{k^{+}}{\Lambda}=\log\frac{\vee}{\Lambda}-\log\frac{\vee}{k^{+}}$. The first two terms are exactly your two log terms, with $(x-z)\cdot(z-x)/[(x-z)^{2}(z-x)^{2}]=-1/(x-z)^{2}$ used.

\section{Checks}
\begin{enumerate}
\item \textbf{Nothing is counted twice.} $N_{2b\,\rm UV}$ is $+\int S$ and the finite and log parts contain $-\int S$: the subtraction line of \eqref{eq:finite} plus the $-2\log$ term of \eqref{eq:logs}. Their sum is the original \eqref{eq:N2}.
\item \textbf{$N_{2\,\rm finite}$ has no UV log.} At $x\to z$, the main terms of \eqref{eq:finite} give $\int_{0}^{1}d\xi\,\big[\xi(1-\xi)+[\tfrac{1}{1-\xi}]_{+}+[\tfrac{1}{\xi}]_{+}-2\big]K_{S}\Cz$, which the subtraction line cancels term by term at each $p^{+}$. The $\mathcal{C}^{(a)}$ part is finite by colour (Step 3).
\item \textbf{$N_{2\,\log}$ has no UV log.} At $x\to z$ the three terms of \eqref{eq:logs} go as $(+1+1-2)\,\frac{(x'-w')\cdot(y-z)}{(x'-w')^{2}(y-z)^{2}(x-z)^{2}}\Cz=0$, and the $\mathcal{C}^{(a)}$ part vanishes by colour.
\item \textbf{The only $\log\mu$ is in $N_{2b\,\rm UV}$.} It is $\big[2\log\frac{k^{+}}{\Lambda}-\frac{11}{6}\big]\times\UVb$. If the $x$ integral of the $-2\log$ term in \eqref{eq:logs} is undone, it is $-2\log\frac{k^{+}}{\Lambda}\,K_{S}\Cz$, which cancels the $+2\log\frac{k^{+}}{\Lambda}$ of $N_{2b\,\rm UV}$ (the $\pm2N_{c}$ blocks in the JIMWLK file). What remains of the UV part is $-\frac{11}{6}\times\UVb$, which goes to the Divergent file and combines with Row II.
\item \textbf{Large $x$ in the log terms.} At $p^{+}=0$ the denominator no longer grows with $x$, and the first term of \eqref{eq:logs} falls only like $1/(2x^{2})$. This is large-distance (IR) behaviour, not UV, and has to be checked in the sum of all $\log(k^{+}/\Lambda)$ terms in the JIMWLK file.
\end{enumerate}

\end{document}
